% =====================================================================
%  패키지 및 서식 설정
% =====================================================================
% ---- 한글 (XeLaTeX + kotex) ----
\usepackage{kotex}
\setmainhangulfont{NanumMyeongjo}[BoldFont=NanumMyeongjoBold,ItalicFont=NanumMyeongjo]
\setsanshangulfont{NanumGothic}[BoldFont=NanumGothicBold,ItalicFont=NanumGothic]
\setmonohangulfont{NanumGothicCoding}

% ---- 쪽 여백 / 줄간격 ----
\usepackage[top=30mm,bottom=30mm,left=30mm,right=25mm,marginparwidth=20mm]{geometry}
\usepackage{setspace}
\onehalfspacing

% ---- 수식 / 정리 ----
\usepackage{amsmath,amssymb,amsthm,mathtools,bm}
\newtheorem{theorem}{정리}[chapter]
\newtheorem{lemma}[theorem]{보조정리}
\newtheorem{proposition}[theorem]{명제}
\newtheorem{corollary}[theorem]{따름정리}
\theoremstyle{definition}
\newtheorem{definition}[theorem]{정의}
\newtheorem{example}[theorem]{예제}
\theoremstyle{remark}
\newtheorem{remark}[theorem]{비고}
\renewcommand{\proofname}{증명}

% ---- 그림 / 표 ----
\usepackage{graphicx}
\graphicspath{{figures/}}
\usepackage{booktabs,multirow,array,tabularx}
\usepackage{caption,subcaption}
\usepackage{float}

% ---- 알고리즘 / 코드 ----
\usepackage[ruled,vlined,linesnumbered]{algorithm2e}
\SetAlgorithmName{알고리즘}{알고리즘}{알고리즘 목록}
\SetKwInput{KwIn}{입력}
\SetKwInput{KwOut}{출력}
\usepackage{listings}
\lstset{basicstyle=\ttfamily\small,breaklines=true,frame=single,
        numbers=left,numberstyle=\tiny,columns=fullflexible}

% ---- 기타 ----
\usepackage{xcolor}
\usepackage{enumitem}
\usepackage{siunitx}

% ---- 한글 캡션 / 목차 이름 ----
\renewcommand{\contentsname}{목차}
\renewcommand{\listfigurename}{그림 목차}
\renewcommand{\listtablename}{표 목차}
\renewcommand{\figurename}{그림}
\renewcommand{\tablename}{표}
\renewcommand{\appendixname}{부록}
\renewcommand{\chaptername}{제}
\usepackage{titlesec}
\titleformat{\chapter}[display]
  {\normalfont\huge\bfseries\centering}{제 \thechapter\ 장}{12pt}{\Huge}
\titlespacing*{\chapter}{0pt}{20pt}{30pt}

% ---- 참고문헌 (biblatex + biber) ----
\usepackage{csquotes}
\usepackage[backend=biber,style=ieee,sorting=nyt]{biblatex}
\addbibresource{references.bib}

% ---- 초안 도구: 할 일 메모 ----
%   \todo{...} 로 메모를 남기면 여백에 표시됩니다.
%   최종본에서는 \documentclass 뒤에 \def\FinalVersion{} 를 넣으면 숨겨집니다.
\usepackage{todonotes}
\ifdefined\FinalVersion
  \presetkeys{todonotes}{disable}{}
\fi

% ---- 하이퍼링크 (마지막에 로드) ----
\usepackage[hidelinks,unicode]{hyperref}
\hypersetup{pdftitle={\ThesisTitleKo},pdfauthor={\AuthorKo}}
\usepackage[nameinlink,noabbrev]{cleveref}
\crefname{figure}{그림}{그림}
\crefname{table}{표}{표}
\crefname{equation}{식}{식}
\crefformat{chapter}{#2제#1장#3}
\Crefformat{chapter}{#2제#1장#3}
\crefformat{section}{#2#1절#3}
\Crefformat{section}{#2#1절#3}
\crefname{theorem}{정리}{정리}
\crefname{definition}{정의}{정의}
\crefname{lemma}{보조정리}{보조정리}
